---
title: "The Wiener Chaos Decomposition, Explained Like a Fourier Series"
author: "Eugen Nesbakken"
date: "2026-07-30"
description: "Wiener chaos decomposes random variables on Gaussian space by polynomial degree, the same way Fourier series decomposes functions by frequency."
category: mathematics
tags: [probability, malliavin-calculus, functional-analysis]
---

\section{Introduction}

I've been working through Malliavin calculus lately, and the first real construction you hit, before any derivatives show up, is the Wiener chaos decomposition. When I first saw it, it looked like a pile of unmotivated definitions: Hermite polynomials, isonormal Gaussian processes, closed subspaces $\mathcal H_n$. But once it clicked, I realized I already knew this story, just in a different costume. It's a Fourier series.

Recall what a Fourier series does to a function on the circle:
$$
f = \sum_{n \in \mathbb{Z}} \hat f(n)\, e^{in\theta}, \qquad \|f\|_{L^2}^2 = \sum_n |\hat f(n)|^2.
$$
You take a hard, arbitrary object, an $L^2$ function, and split it into orthogonal pieces indexed by frequency $n$. Different frequencies never interact, and the total energy of $f$ is just the sum of the energy in each frequency. That's it, that's the whole trick, and it's why Fourier series are useful everywhere: you've turned one hard infinite-dimensional problem into countably many easy one-dimensional problems.

Wiener chaos does exactly this, except for random variables built out of a Gaussian process, and the role of "frequency" is played by \emph{polynomial degree}. That's the whole post: I want to walk through the construction and show you where every piece of the Fourier picture has a chaos-theoretic twin.

\section{Setup: an isonormal Gaussian process}

First we need the Gaussian analogue of "functions on the circle." Let $H$ be a real separable Hilbert space with inner product $\langle \cdot, \cdot \rangle_H$.

\begin{definition}[Isonormal Gaussian process]
A stochastic process $W = \{W(h), h \in H\}$ on a complete probability space $(\Omega, \mathcal F, P)$ is an \emph{isonormal Gaussian process} if $W$ is a centered Gaussian family with
$$
E[W(h)W(g)] = \langle h, g \rangle_H \qquad \text{for all } h,g \in H.
$$
\end{definition}

Don't overthink this: $H$ is just the space of "directions" you're allowed to feed in, and $W(h)$ is the Gaussian random variable you get back. If $H = L^2([0,\infty))$ and $W$ is Brownian motion, then $W(h) = \int_0^\infty h(t)\,dW_t$. That's nothing but the Wiener integral. The inner product structure on $H$ becomes the covariance structure of the Gaussian family. That's the entire content of the definition.

We write $\mathcal G$ for the $\sigma$-field generated by $\{W(h), h \in H\}$, and everything below happens inside $L^2(\Omega, \mathcal G, P)$.

\section{Hermite polynomials: the right basis}

On the circle, the natural orthogonal basis is $\{e^{in\theta}\}$. On Gaussian space, it turns out the natural orthogonal basis is built from Hermite polynomials, and once you see why, the rest of the story falls into place.

Define $H_n$ by the Rodrigues formula
$$
H_n(x) = \frac{(-1)^n}{n!} e^{x^2/2} \frac{d^n}{dx^n} e^{-x^2/2}, \qquad n \ge 1, \qquad H_0 = 1.
$$
You don't need to stare at that formula for long, because there's a much better way to get your hands on these: they are the coefficients of a generating function,
$$
F(x,t) = \exp\!\left(tx - \frac{t^2}{2}\right) = \sum_{n=0}^\infty t^n H_n(x).
$$
You can see this yourself by completing the square, $tx - t^2/2 = \tfrac12 x^2 - \tfrac12(x-t)^2$, and then Taylor-expanding the resulting expression in $t$ around $t = 0$. The $n$-th derivative in $t$ at $t=0$, divided by $n!$, hands you exactly $H_n(x)$. The first two are $H_1(x) = x$ and $H_2(x) = \tfrac12(x^2 - 1)$, and they satisfy the recursion
$$
(n+1)H_{n+1}(x) = xH_n(x) - H_{n-1}(x).
$$

Here's the fact that actually does all the work, and it's the direct analogue of $\int_0^{2\pi} e^{in\theta} e^{-im\theta}\, d\theta = 0$ for $n \ne m$:

\begin{lemma}
Let $X, Y$ be jointly Gaussian with $E(X) = E(Y) = 0$ and $E(X^2) = E(Y^2) = 1$. Then for all $n, m \ge 0$,
$$
E[H_n(X)H_m(Y)] =
\begin{cases}
0 & n \ne m, \\[4pt]
\dfrac{1}{n!}\big(E[XY]\big)^n & n = m.
\end{cases}
$$
\end{lemma}

I like this proof because it's genuinely just a trick, not hard analysis. Plug $sX - s^2/2$ and $tY - t^2/2$ into the generating function, multiply the two together, and take expectations:
$$
E\!\left[\exp\!\left(sX - \tfrac{s^2}{2}\right)\exp\!\left(tY - \tfrac{t^2}{2}\right)\right].
$$
The exponent simplifies to $sX + tY - (s^2+t^2)/2$, and $Z := sX + tY$ is itself a centered Gaussian, with variance $s^2 + 2st\,E[XY] + t^2$. Using $E[e^Z] = e^{\sigma_Z^2/2}$, everything collapses to
$$
E\!\left[\exp\!\left(sX - \tfrac{s^2}{2}\right)\exp\!\left(tY - \tfrac{t^2}{2}\right)\right] = \exp\big(st\,E[XY]\big).
$$
Now just read off the coefficient of $s^nt^m$ on both sides. On the left it's $E[H_n(X)H_m(Y)]$ by definition of the generating function, and on the right, $\exp(st\,E[XY]) = \sum_n \frac{(st\,E[XY])^n}{n!}$ only has terms where the power of $s$ matches the power of $t$. That forces $n = m$, and matching coefficients gives exactly the lemma.

So no matter how correlated $X$ and $Y$ are, different Hermite degrees stay orthogonal. Degree is doing the job that frequency does for the circle: it's a label that automatically kills cross-terms.

\section{The chaos spaces}

Now define, for each $n \ge 1$, $\mathcal H_n$ to be the closed linear span in $L^2(\Omega, \mathcal F, P)$ of
$$
\{H_n(W(h)) : h \in H,\ \|h\|_H = 1\},
$$
with $\mathcal H_0 = \mathbb{R}$ (just the constants). This is the \emph{Wiener chaos of order $n$}. The lemma above hands us orthogonality for free: $\mathcal H_n \perp \mathcal H_m$ for $n \ne m$, since $H_n(W(h))$ and $H_m(W(h'))$ have zero covariance regardless of how $h$ and $h'$ are correlated.

That's the easy half. The interesting question, same as always, is whether these orthogonal pieces actually cover everything: the analogue of "trigonometric polynomials are dense in $L^2$ of the circle."

\begin{theorem}[Chaos decomposition]
\label{chaos-decomp}
$$
L^2(\Omega, \mathcal G, P) = \bigoplus_{n=0}^\infty \mathcal H_n.
$$
\end{theorem}

The proof strategy is the same one you'd use to prove completeness of any orthogonal system: show that anything orthogonal to every single piece has to be zero. Suppose $X \perp \mathcal H_n$ for every $n$. Since $x^n$ is a linear combination of $H_0(x), \dots, H_n(x)$, this gives $E[X\,W(h)^n] = 0$ for every $n \ge 0$ and every unit vector $h$, and summing the Taylor series then gives
$$
E\big[X\exp(tW(h))\big] = 0 \qquad \text{for all } t \in \mathbb{R},\ h \in H.
$$
The last ingredient is that random variables of the form $e^{W(h)}$ are \emph{total} in $L^2(\Omega, \mathcal G, P)$: their linear span is dense, so having zero inner product with all of them forces $X = 0$. (This is the one step I had to just accept on faith the first time through; the density argument itself is a bit more involved than I want to get into here.)

\section{Projections, and the picture to keep in your head}

Once you have the decomposition, you get an orthogonal projection $J_n : L^2(\Omega) \to \mathcal H_n$ for each $n$, and every $X \in L^2(\Omega, \mathcal G, P)$ splits as
$$
X = \sum_{n=0}^\infty J_n X, \qquad \|X\|_{L^2}^2 = \sum_{n=0}^\infty \|J_n X\|_{L^2}^2.
$$
This second identity is just Parseval again, and I think the cleanest way to hold onto it is geometrically: the chaoses $\mathcal H_0, \mathcal H_1, \dots$ are mutually perpendicular axes in $L^2(\Omega)$ under the inner product $\langle X, Y \rangle = E[XY]$, and $\|X\|^2 = \sum_n \|J_nX\|^2$ is literally the Pythagorean theorem applied to $X = \sum_n J_nX$.

\includegraphics[width=0.75\textwidth]{images/tikz/wiener-chaos-axes.png}

I find it genuinely useful to picture $X$ as a single vector sitting in this infinite-dimensional space, with $J_0X, J_1X, J_2X, \dots$ its shadow on each axis: the constant part, the "linear-in-$W$" part, the "quadratic-in-$W$" part, and so on, all orthogonal to each other, all adding up in quadrature to give back $\|X\|^2$.

One thing worth flagging, because it tripped me up: it's tempting to think each chaos is built purely out of finite-degree polynomials in finitely many $W(h_i)$'s, full stop. Let $\mathcal P_n^0$ be the space of variables $p(W(h_1), \dots, W(h_k))$ with $\deg p \le n$, not yet closed, and $\mathcal P_n$ its $L^2$-closure. It's true that $\mathcal P_n = \bigoplus_{i=0}^n \mathcal H_i$, but the closure is not just a technicality; it's doing real work.

\begin{example}[the closure is strictly bigger than $\mathcal P_n^0$]
Let $e_1, e_2, \dots$ be orthonormal in $H$, so that $W(e_1), W(e_2), \dots$ are i.i.d.\ standard Gaussians, and pick $\sum_i c_i^2 < \infty$ with infinitely many $c_i \ne 0$. The partial sums $X_N = \sum_{i \le N} c_i W(e_i)$ live in $\mathcal P_1^0$, and if $X = \sum_i c_i W(e_i)$, then $E[(X-X_N)^2] = \sum_{i>N} c_i^2 \to 0$, so $X \in \mathcal P_1$. But $X$ genuinely depends on infinitely many directions at once, so it can't be written as $p(W(h_1), \dots, W(h_k))$ for any single finite family. So $X \in \mathcal P_1 \setminus \mathcal P_1^0$.
\end{example}

This is exactly the same phenomenon as an $L^2$ function on the circle whose Fourier series converges but which isn't itself a finite trigonometric polynomial. Closing the space is precisely what lets you admit these infinite, convergent superpositions.

\section{An explicit basis, and the case that makes it all click}

Higher-dimensional Fourier analysis (say, on a torus) uses products $e^{in_1\theta_1}\cdots e^{in_k\theta_k}$ indexed by multi-indices. Exactly the same move works here. Fix an orthonormal basis $\{e_i\}_{i\ge1}$ of $H$, let $\Lambda$ be the set of multi-indices $a = (a_1, a_2, \dots)$ with only finitely many nonzero entries, and set
$$
\Phi_a = \sqrt{a!} \prod_{i=1}^\infty H_{a_i}(W(e_i)), \qquad a! := \prod_i a_i!.
$$
This is a finite product (since $H_0 = 1$), and using the orthogonality lemma plus independence of $W(e_1), W(e_2), \dots$, you can check directly that $\{\Phi_a : a \in \Lambda\}$ is orthonormal. Better still, $\{\Phi_a : |a| = n\}$, where $|a| = \sum_i a_i$, is a \emph{complete} orthonormal basis of $\mathcal H_n$. Total degree plays the role of total frequency on the torus.

\includegraphics[width=0.85\textwidth]{images/tikz/fourier-chaos-dictionary.png}

And here's the case that made the whole thing click for me: the simplest possible one.

\begin{example}[$H = \mathbb{R}$]
Take $(\Omega, \mathcal F, P) = (\mathbb{R}, \mathcal B(\mathbb{R}), \nu)$ with $\nu = N(0,1)$, and $H = \mathbb{R}$ with $W(h)(x) = hx$. The only unit vectors are $h = \pm1$. Here $\mathcal H_n$ is one-dimensional, spanned by $H_n(x)$, and Theorem \ref{chaos-decomp} says exactly this: the Hermite polynomials $\{H_n\}_{n\ge0}$ form a complete orthonormal system in $L^2(\mathbb{R}, \nu)$.
\end{example}

That's a fact I already knew, Hermite polynomials as an orthogonal basis of $L^2$ of a Gaussian measure, and seeing it fall out of the general chaos decomposition as the one-dimensional case is what convinced me the analogy isn't just decoration. It's the same statement, one variable at a time, the way a single Fourier coefficient sits inside the full basis on a torus.

\begin{remark}
Where this is actually going: the chaos decomposition is the launching pad for the Malliavin derivative $D$ and its adjoint $\delta$, and they act on each $\mathcal H_n$ almost diagonally, the same way differentiation multiplies the $n$-th Fourier mode by $in$. That's next, once I've actually finished reading the proof I'm currently stuck on.
\end{remark}
